\section{Introduction}
$\infty$-operads are combinatorial gadgets for controlling homotopy-coherent algebraic theories. For example, there exist $\infty-$operads $E_n$ whose algebras are exactly $E_n$-algebras. Lurie defines (symmetric) operads as certain inner fibrations over $\mathbb{F}_*$, the category of pointed finite sets \cite{HA}, which should be thought of as a generalisation of the category of operators associated to a $1-$operad. Lurie's definition is known to be equivalent other models of $\infty-$operads, namely the dendroidal $\infty-$operads of Moerdijk-Weiss \cite{Moerdijk-Weiss} \cite{Heuts-Hinich-Moerdijk} \cite{Hinich-Moerdijk}, and the complete Segal $\infty-$operads of Barwick \cite{Bar18} \cite{Chu-Haugseng-Heuts}.

Some algebras may be better described using a different formulation entirely. For example, the associative operad $Assoc$ is more easily described as the terminal \defn{non-symmetric} operad, i.e. as a certain fibration over $\Delta^{op}$. Barwick's notion of a perfect operator category \cite{Bar18} provides sufficient conditions for a category $\phi$ to admit a sensible notion of $\phi-$operads. This is further generalised by the algebraic patterns of Chu and Haugseng \cite{ChuHaugseng}. Noting that symmetric monoidal $\infty-$categories are a special case of $\infty-$operads, these formulations also generalise the notion of symmetric monoidal category.

A structure not encaptured by either notion is that of a duoidal category. This is, roughly, a category $C$ equipped with two symmetric tensor products $(\tensor_A,1_A)$, $(\tensor_B,1_B)$ such that one distributes laxly over the other. The 1-categorical version is due to Aguiar and Mahajan \cite{Aguiar-Mahajan}, under the name of 2-monoidal category, and an $\infty-$categorical generalisation has been developed by \cite{Torii24} in the form of certain mixed fibrations to $\Delta^{op}.$ When the two units for the tensor products coincide, we call the category twofold-monoidal. Twofold monoidal $1-$categories have been studied by \cite{Balteanu-Fiedorowicz-Schwaenzl-Vogt} and \cite{BeilinsonDrinfeld}, in the latter under the name of compound tensor category. An $\infty-$categorical version is proposed by Barwick \cite{Bar24}, as certain fibrations over $\E_*$, the category of pointed reflexive cographs (see \cref{sec:cograph}).

We are interested in describing an operadic version of this concept, giving an $\infty-$categorical version of compound pseudo-tensor categories in the sense of \cite{BeilinsonDrinfeld}.

The paper is organised as follows. Chapter $1$ is devoted to a literature review. In \cref{sec:operator-category} we will review the notion of an operator category $\phi$ and the related notion of $\phi-$operad. \cref{sec:envelope} is devoted to the symmetric monoidal envelope, and the proof by Haugseng and Kock \cite{HaugsengKock} that operads form a full subcategory of the category of symmetric monoidal functors equipped with a symmetric monoidal functor to $\mathbb{F}^{\sqcup}$. Our purpose of reviewing this result is to speculate on the existence of a twofold monoidal envelope comonad on the category of twofold monoidal categories, from which we can \textit{define} twofold operads as certain twofold monoidal categories with a twofold monoidal functor to $\E$. 

%Operads are closely related to the notion of symmetric monoidal category. There is an obvious forgetful functor from the category of operads to the category of symmetric monoidal categories which has a left adjoint given by the symmetric monoidal envelope \cite{HA}. By a theorem of Haugseng and Kock \cite{HaugsengKock}, this left adjoint is fully faithful when considered as a functor to symmetric monoidal categories equipped with a symmetric monoidal functor to $\mathbb{F}^{\sqcup}$, finite sets with the monoidal product given by disjoint union. Our purpose of reviewing this result is to speculate on the existence of a twofold monoidal envelope comonad on the category of twofold monoidal categories, from which we can \textit{define} twofold operads as certain twofold monoidal categories with a monoidal functor to $\E$ with its canonical twofold monoidal structure. 

Chapter 2 is devoted to twofold monoidal structures. In \cref{sec:cograph} we describe cographs and the categories $\E$ and $\mathbb{D}$, and in \cref{sec:twofold-monoidal}, review Barwick's definition of a twofold monoidal category. $\mathbb{D}$ and $\E$ are free twofold monoidal on a monoid for, respectively, the left and the right product. In \cref{sec:H} we introduce the category $\mathbb{H}$, which is the free twofold monoidal category on a monoid for the left product that is a comonoid for the right product. We speculate on a definition of twofold operads described as a fibration over $\Lambda(\mathbb{H}).$ Finally, in \cref{sec:whats-next} we discuss open questions and potential next steps for the project.

\subsection{Conventions}
\begin{itemize}
\item We use implicit $\infty$-conventions. A category will always refer to an $(\infty,1)$-category, an operad to an $(\infty,1)-$operad, and so on, unless explicitly specified.
\item An operad will always refer to the coloured or multicategory case, unless explicitly specified. This is, broadly speaking, in contrast to the language used in the 1-categorical literature, but consistent with the language in the $\infty-$categorical literature.
\item A symmetric monoidal functor will mean \textit{strong} symmetric monoidal functor, i.e. an edge of $Cat_{\infty/\mathbb{F}_*}$ between symmetric monoidal categories that preserves coCartesian edges. Since we only consider commutative structures in this paper, we will sometimes leave "symmetric" implicit and so talk about monoidal functors of monoidal categories.
\end{itemize}